% GENERATED FILE. Edit canonical lessons, then run npm run generate:books.
% canonical-source-hash: fnv1a64:8c01aa4d828a9c49
% canonical-lessons: GU-C323-sahid, GU-C323-balidan, GU-C323-vijay, GU-C323-ver, GU-C323-padkarvu

\chapter{The Martyrs}
\label{ch:gu-the-martyrs}
\hlchaptermodality{323}

\begin{chapteropening}
\textbf{By the end of this chapter:} \emph{I can say who gave their lives and how they are remembered.}

Complete the last lesson of chapter 323: I can say who gave their lives and how they are remembered.
\end{chapteropening}

\section[\texorpdfstring{\gu{શહીદ} (śahīd)}{શહીદ (śahīd)}]{\gu{શહીદ} (śahīd) — a martyr}

\label{lesson:GU-C323-sahid}



\begin{quote}
Before the new one: say the Gujarati for to slay, then the Gujarati for to fight on.
\end{quote}

\subsection*{You'll want to know: \gu{શહીદ}}
\textbf{\gu{શહીદ}} — \emph{śahīd} — \textquotedblleft{}a martyr\textquotedblright{}.

\begin{quote}
\gu{ગામ} \gu{આ} \gu{શહીદોનું} \gu{ઋણી} \gu{છે}.
\end{quote}

\begin{quote}
\emph{gām ā śahīdonũ ṛṇī che.}
\end{quote}

\begin{quote}
The village owes a debt to these martyrs.
\end{quote}

A masculine or feminine noun, by the person meant.

\begin{cousinweb}
\textbf{\gu{શહીદ}} (\emph{śahīd}) is usually traced, through Persian, to Arabic \emph{šahīd}, \textquotedblleft{}a witness\textquotedblright{}, and so \textquotedblleft{}a martyr\textquotedblright{}, one whose death bears witness. Someone \textbf{\gu{શહીદ} \gu{થાય} \gu{છે}}, dies for a cause.
\end{cousinweb}

\begin{grammarlens}[title={What you've built}]
One more word for telling what happened.
\end{grammarlens}

\subsection*{Your turn}
\begin{itemize}
\raggedright
  \item \emph{Say it:} \emph{śahīd}
  \item \emph{Say it:} \emph{śahīd}, once more
  \item \emph{Say it:} \emph{gām ā śahīdonũ ṛṇī che.} — The village owes a debt to these martyrs.
  \item \emph{Recall:} say the Gujarati for to slay, then the Gujarati for to fight on, then say \emph{śahīd} again
\end{itemize}

\begin{tcolorbox}[breakable,colback=teal!4,colframe=teal!35!black,title={Before you move on}]
What does \gu{શહીદ} mean? (\textquotedblleft{}A martyr\textquotedblright{}.) Say it once more.
\end{tcolorbox}
\section[\texorpdfstring{\gu{બલિદાન} (balidān)}{બલિદાન (balidān)}]{\gu{બલિદાન} (balidān) — a sacrifice}

\label{lesson:GU-C323-balidan}



\begin{quote}
Before the new one: say the Gujarati for to fight on, then the Gujarati for a martyr.
\end{quote}

\subsection*{You'll want to know: \gu{બલિદાન}}
\textbf{\gu{બલિદાન}} — \emph{balidān} — \textquotedblleft{}a sacrifice\textquotedblright{}.

\begin{quote}
\gu{સૈનિકોના} \gu{બલિદાનથી} \gu{ગામ} \gu{બચી} \gu{ગયું}.
\end{quote}

\begin{quote}
\emph{sainikonā balidānthī gām bacī gayũ.}
\end{quote}

\begin{quote}
The soldiers' sacrifice saved the village.
\end{quote}

A neuter noun.

\begin{cousinweb}
\textbf{\gu{બલિદાન}} (\emph{balidān}) is Sanskrit \emph{balidāna}, \textquotedblleft{}the presenting of an offering\textquotedblright{} (Monier-Williams), built from \emph{bali}, an offering, and \emph{dāna}, a giving. In Gujarati it is giving up one's life or comfort for others.
\end{cousinweb}

\begin{grammarlens}[title={What you've built}]
One more word for telling what happened.
\end{grammarlens}

\subsection*{Your turn}
\begin{itemize}
\raggedright
  \item \emph{Say it:} \emph{balidān}
  \item \emph{Say it:} \emph{balidān}, once more
  \item \emph{Say it:} \emph{sainikonā balidānthī gām bacī gayũ.} — The soldiers' sacrifice saved the village.
  \item \emph{Recall:} say the Gujarati for to fight on, then the Gujarati for a martyr, then say \emph{balidān} again
\end{itemize}

\begin{tcolorbox}[breakable,colback=teal!4,colframe=teal!35!black,title={Before you move on}]
What does \gu{બલિદાન} mean? (\textquotedblleft{}A sacrifice\textquotedblright{}.) Say it once more.
\end{tcolorbox}
\section[\texorpdfstring{\gu{વિજય} (vijay)}{વિજય (vijay)}]{\gu{વિજય} (vijay) — a victory}

\label{lesson:GU-C323-vijay}



\begin{quote}
Before the new one: say the Gujarati for a martyr, then the Gujarati for a sacrifice.
\end{quote}

\subsection*{You'll want to know: \gu{વિજય}}
\textbf{\gu{વિજય}} — \emph{vijay} — \textquotedblleft{}a victory\textquotedblright{}.

\begin{quote}
\gu{છેવટે} \gu{રાજાની} \gu{સેનાનો} \gu{વિજય} \gu{થયો}.
\end{quote}

\begin{quote}
\emph{chevṭe rājānī senāno vijay thayo.}
\end{quote}

\begin{quote}
In the end the king's army won.
\end{quote}

A masculine noun.

\begin{cousinweb}
\textbf{\gu{વિજય}} (\emph{vijay}) is a tatsama, Sanskrit \emph{vijaya}, \textquotedblleft{}victory, conquest\textquotedblright{} (Monier-Williams). It happens with \textbf{\gu{થવું}}: \textbf{\gu{સેનાનો} \gu{વિજય} \gu{થયો}}, the army won.
\end{cousinweb}

\begin{grammarlens}[title={What you've built}]
One more word for telling what happened.
\end{grammarlens}

\subsection*{Your turn}
\begin{itemize}
\raggedright
  \item \emph{Say it:} \emph{vijay}
  \item \emph{Say it:} \emph{vijay}, once more
  \item \emph{Say it:} \emph{chevṭe rājānī senāno vijay thayo.} — In the end the king's army won.
  \item \emph{Recall:} say the Gujarati for a martyr, then the Gujarati for a sacrifice, then say \emph{vijay} again
\end{itemize}

\begin{tcolorbox}[breakable,colback=teal!4,colframe=teal!35!black,title={Before you move on}]
What does \gu{વિજય} mean? (\textquotedblleft{}A victory\textquotedblright{}.) Say it once more.
\end{tcolorbox}
\section[\texorpdfstring{\gu{વેર} (ver)}{વેર (ver)}]{\gu{વેર} (ver) — enmity; a feud}

\label{lesson:GU-C323-ver}



\begin{quote}
Before the new one: say the Gujarati for a sacrifice, then the Gujarati for a victory.
\end{quote}

\subsection*{You'll want to know: \gu{વેર}}
\textbf{\gu{વેર}} — \emph{ver} — \textquotedblleft{}enmity; a feud\textquotedblright{}.

\begin{quote}
\gu{બંને} \gu{રાજાઓ} \gu{વચ્ચે} \gu{જૂનું} \gu{વેર} \gu{હતું}.
\end{quote}

\begin{quote}
\emph{banne rājāo vacce jūnũ ver hatũ.}
\end{quote}

\begin{quote}
There was an old feud between the two kings.
\end{quote}

A neuter noun.

\begin{cousinweb}
\textbf{\gu{વેર}} (\emph{ver}) is inherited from Sanskrit \emph{vaíra}, \textquotedblleft{}enmity\textquotedblright{}, and CDIAL lists the Gujarati noun, \textquotedblleft{}enmity, revenge\textquotedblright{} (CDIAL 12142). It is neuter: \textbf{\gu{જૂનું} \gu{વેર}}, an old feud.
\end{cousinweb}

\begin{grammarlens}[title={What you've built}]
One more word for telling what happened.
\end{grammarlens}

\subsection*{Your turn}
\begin{itemize}
\raggedright
  \item \emph{Say it:} \emph{ver}
  \item \emph{Say it:} \emph{ver}, once more
  \item \emph{Say it:} \emph{banne rājāo vacce jūnũ ver hatũ.} — There was an old feud between the two kings.
  \item \emph{Recall:} say the Gujarati for a sacrifice, then the Gujarati for a victory, then say \emph{ver} again
\end{itemize}

\begin{tcolorbox}[breakable,colback=teal!4,colframe=teal!35!black,title={Before you move on}]
What does \gu{વેર} mean? (\textquotedblleft{}Enmity; a feud\textquotedblright{}.) Say it once more.
\end{tcolorbox}
\section[\texorpdfstring{\gu{પડકારવું} (paḍkārvũ)}{પડકારવું (paḍkārvũ)}]{\gu{પડકારવું} (paḍkārvũ) — to challenge}

\label{lesson:GU-C323-padkarvu}



\begin{quote}
Before the new one: say the Gujarati for a victory, then the Gujarati for enmity.
\end{quote}

\subsection*{You'll want to know: \gu{પડકારવું}}
\textbf{\gu{પડકારવું}} — \emph{paḍkārvũ} — \textquotedblleft{}to challenge\textquotedblright{}.

\begin{quote}
\gu{યોદ્ધાએ} \gu{દુશ્મનને} \gu{પડકાર્યો}.
\end{quote}

\begin{quote}
\emph{yoddhāe duśmanne paḍkāryo.}
\end{quote}

\begin{quote}
The warrior challenged the enemy.
\end{quote}

Past, with \gu{તેણે} (he or she …): \gu{તેણે} \gu{દુશ્મનને} \gu{પડકાર્યો}, \emph{teṇe duśmanne paḍkāryo}.

\begin{cousinweb}
\textbf{\gu{પડકારવું}} (\emph{paḍkārvũ}) No dictionary we use records an older form, so its history is not settled. The one challenged takes \textbf{\gu{ને}}, and the verb agrees with him: \textbf{\gu{દુશ્મનને} \gu{પડકાર્યો}}.
\end{cousinweb}

\begin{grammarlens}[title={What you've built}]
One more word for telling what happened.
\end{grammarlens}

\subsection*{Your turn}
\begin{itemize}
\raggedright
  \item \emph{Say it:} \emph{paḍkārvũ}
  \item \emph{Say it:} \emph{paḍkārvũ}, once more
  \item \emph{Say it:} \emph{yoddhāe duśmanne paḍkāryo.} — The warrior challenged the enemy.
  \item \emph{Recall:} say the Gujarati for a victory, then the Gujarati for enmity, then say \emph{paḍkārvũ} again
\end{itemize}

\begin{tcolorbox}[breakable,colback=teal!4,colframe=teal!35!black,title={Before you move on}]
What does \gu{પડકારવું} mean? (\textquotedblleft{}To challenge\textquotedblright{}.) Say it once more.
\end{tcolorbox}
